\section{问题三：通信约束下的运输与中继调度}

\subsection{双向链路与连续覆盖}

固定网关G01位于O01上方20~m。任意端点$a,b$的双向损耗门限取两个通信方向允许损耗的较小值。若三维距离为$D_{ab}(t)$~km，地形遮挡指示为$B_{ab}(t)$，可用性判据为
\begin{equation}
32.45+20\log_{10}(2400)+20\log_{10}D_{ab}(t)+10B_{ab}(t)\le L^{\max}_{a\leftrightarrow b}.
\end{equation}
由附件收发参数得到三个门限：运输机至网关为122~dB，运输机至中继为116~dB，中继至网关为126~dB。运输机在每一时刻只需直连成立，或同一架中继的接入与回传同时成立。中继悬停点必须在DEM内、离地不超过300~m。

为了验证连续而非仅离散采样的通信，先在架次中的爬升、巡航、下降、交接段以及中继任务切换时刻切分。对子区间以两端点最大三维距离给出自由空间损耗上界；若再加10~dB仍满足门限，则无须依赖视线。其余区间对运输机运动端点与静止通信端点形成的扫掠三角形逐像元检验DEM，对视线高度采用像元全域的保守下界。仍无法证明时继续二分，直到获得可用的直连或同一中继保障证书。这样可覆盖段内全部位置与时刻。

\subsection{联合时间表和资源}

问题二的多点航线已形成至少三个独立服务区组件。先按薄弱轨迹点筛选西、东中继点，再搜索四个北部候选位置、西部两档、东部五档与北部两档服务窗口，共80种组合；对每种组合，在连续通信可行起飞区间内重排运输航线的发车时刻。对较优组合，再尝试把迟到货箱从原架次拆出并改为单箱直送。当前把S002的一箱水改由B型机直送，第三问运输架次因此由24增为25。排程按加权迟到、运输完成时刻依次优化，再对整条路线与全部中继任务进行连续通信回代。此搜索只覆盖有限候选，不给出全局最优性证明。两架中继机在西点W、东点E和北点N共执行5个架次，悬停高度均为对应DEM地面高程上方300~m。R01在W服务两次；R02依次到E、N、E。每一架次计入准备、去程、建链、悬停通信、返航及300~s周转，5组能源组件分别装机，剩余一组备用。图~\ref{fig:relay}显示三个悬停位置，图~\ref{fig:joint}显示运输与中继并行的时间表。

\begin{figure}[H]\centering\includegraphics[width=.86\textwidth]{../figures/relay_map.pdf}\caption{西、东、北中继悬停点及运输区域}\label{fig:relay}\end{figure}
\begin{figure}[H]\centering\includegraphics[width=.90\textwidth]{../figures/joint_gantt.pdf}\caption{问题三运输与中继架次时间表}\label{fig:joint}\end{figure}

\begin{table}[H]\centering\small\caption{通信约束下联合方案的主要结果}\label{tab:q3}
\begin{tabular}{lr}\toprule
指标 & 数值\\\midrule
运输/中继架次 & 25/5\\
运输/中继/总能耗（kWh） & 77.3092/5.8432/83.1524\\
最后运输机/所有设备返回时刻（s） & 10208.43/10801.63\\
优先系数加权迟到 & 1935.45\\
通信证明区间数 & 4173\\
最小链路预算裕量（dB） & 0.0061\\
硬时限违约箱数 & 0\\
\bottomrule\end{tabular}\end{table}

全部80箱唯一交付，最紧硬时限余量112.86~s，最低运输与中继SOC分别为20.1676\%和23.6510\%。通信证书覆盖运输机所有起飞后的阶段及交接区间；两架中继机及其能源组件均无重叠。与问题二相比，运输架次增加1个，中继架次增加5个，总能耗增加7.6107~kWh，全部设备返回时刻晚2118.38~s，加权迟到由零升至1935.45。这个软时效代价不能仅以硬时限零违约概括。最小通信裕量只有0.0061~dB，表明方案满足标称参数但对地形高程和设备损耗扰动敏感。
